\documentclass{article}
\usepackage{amsmath, amssymb}
\usepackage{wasysym}
\usepackage[utf8]{inputenc}
\usepackage{transparent-nometadata, tipa}%fontspec}
\usepackage[legalpaper, portrait, margin=5em, top=100pt]{geometry}
\usepackage{titlesec}
\usepackage{fancyhdr}
\usepackage[none]{hyphenat}
\usepackage{tabularx}
\usepackage{ctable}
\usepackage{array}
\usepackage{longtable, needspace}
%\usepackage{pgffor}
\usepackage{multido, etoolbox}
\usepackage{xparse, xstring}
\usepackage{cancel}
\usepackage{expl3}
%\usepackage{unicode-math}
%\newfontfamily\ipafont{Times New Roman}
\newcommand{\ipafont}[1]{\textipa{#1}}
\newcommand{\tab}{\hspace{4em}}
\newcommand{\myname}{D\hat{\jmath}_\prime mz}

%\Umathchardef\xnot="3 \symoperators "0338
%\AtBeginDocument{
%  \renewcommand\not[1]{#1\mathrel{\mkern1mu}\xnot}
%  \renewcommand{\notin}{\not\in}
%  \renewcommand{\notni}{\not\ni}
%}

\ExplSyntaxOn
\cs_new:Npn \retokenize #1 #2 {
  \tl_set_rescan:Nnn #1 {} {#2}
}
\ExplSyntaxOff
\NewDocumentCommand{\arraylen}{m}{%
  \setcounter{arraylength}{0}%
  \@for\item:=#1\do{%
    \stepcounter{arraylength}%
  }%
  \thearraylength
}

\titleformat{\section}{\fontsize{24pt}{18pt}\selectfont\bfseries}{\thesection}{1em}{}

\renewcommand{\arraystretch}{1.5}

\newcolumntype{M}[1]{>{\centering\arraybackslash}m{#1}}
\newcolumntype{?}{!{\vrule width 3pt}}

\newcommand{\tbs}{\textbackslash}

\newcommand{\mataHello}{\tilde{\bar{E}}\hat{y}}%{{\tilde{E}\mkern-10mu\raise1pt\hbox{$/\mkern-8mu/$}}\mkern2mu\hat{y}}

\NewDocumentCommand{\getcolcount}{ m O{ampcount}}{%
	\typeout{received: #1}
	\begingroup
	\catcode`\&=12
	\def\tmpstring{#1}%
	\scantokens\expandafter{\tmpstring}%
	\StrCount{\tmpstring}{&}[\tmpcount]%
	\typeout{counted: \tmpcount}
	\xdef\tmpcount{\the\numexpr\tmpcount+1\relax}%
	\endgroup
	\expandafter\edef\csname#2\endcsname{\tmpcount}
}

%\newcommand\wttitlespec{English & Mata}%
\NewDocumentEnvironment{wordtabular}{ O{4em} O{3em} m }
{%
\def\colspec{}%
\def\titlespec{English & Mata}
\multido{}{#3}{%
	\gappto\colspec{||M{#1}|M{#2}}%
}
\ifnum#3>1
	\multido{}{\numexpr#3-1\relax}{%
		\gappto\titlespec{ & English & Mata}
	}
\fi
\gappto\colspec{||}
\expandafter\tabular\expandafter{\colspec}
\specialrule{.2em}{1em}{0em}\hline
\titlespec \\ \hline}
{\endtabular}

\NewDocumentEnvironment{stdtable}{ m }
{%
\def\colspec{|}%
\multido{}{#1}{%
	\gappto\colspec{|c}
}%
\gappto\colspec{||}
\expandafter\tabular\expandafter{\colspec}
\specialrule{.2em}{1em}{0em}%\hline
}
{\endtabular}

\makeatletter
\newcommand{\headline}[1]{\gdef\@headline{#1}}
\newcommand{\subheadline}[1]{\gdef\@subheadline{#1}}
\renewcommand{\maketitle}{
  \begin{center}
  \topskip0pt
  \vspace*{20em}
    {\fontsize{144pt}{28pt}\selectfont \@title \par}
    \vspace{2em}
    {\fontsize{36pt}{16pt}\selectfont \@author \par}
    \vspace{5em}
    {\fontsize{36pt}{16pt}\selectfont \textbf \@headline \par}
    \vspace{2em}
    {\fontsize{24pt}{16pt}\selectfont \@subheadline \par}
    \vspace{\fill}
    {\fontsize{18pt}{14pt}\selectfont \@date \par}
  \vspace*{3em}
  \end{center}
}
\makeatother


\title{Mata}
\author{J. Mann}
\headline{A comprehensive guide to the language}
\subheadline{Written in English}
\date{April 2026}

\begin{document}
%\fontsize{20pt}{24pt}\selectfont
\begin{titlepage}
\maketitle
\end{titlepage}
\newpage
\pagestyle{fancy}
\lhead{\fontsize{36pt}{12pt}\selectfont Mata}
\rhead{}

\setcounter{section}{-1}
\section{Introduction}
\vspace{-0.5em}
{\fontsize{16pt}{12pt}\selectfont
$\mathbf{\mataHello}${\fontsize{8pt}{12pt}\selectfont(Ehh-ih-ya)} $\mathbf{|\ Hello}$
\vspace{0.2em}\\
Welcome to this comprehensive guide to my language known as \emph{Mata}.\\
This is my first conlang, made to revolve around mathematical notation, in particular set notation.
The core concept to know is that almost everything can be described as a set, containing all of the information about said thing. For example, a person can stated to be a set (i.e. $\mathbb{P}$) that contains all of the adjectives describing them. We can then make statements about that set (i.e "$\mathbb{P}\ni\text{stupid}$" \emph{aka} the person is stupid, however "stupid" wouldn't just be in English).
}
\newpage

\vspace{10em}
\section{Lexicon}
\subsection{Alphabet}
Using the IPA for the sounds\\
\begin{tabular}{||c|c||c|c||c|c||c|c||c|c||}
\specialrule{.2em}{1em}{0em}
\hline
Symbol & Sound & Symbol & Sound & Symbol & Sound & Symbol & Sound & Symbol & Sound \\
\hline
Aa & æ & Gg & g & Mm & m & Ss & s & Yy & j \\
\hline
Bb & b & Hh & h & Nn & n & Tt & t & Zz & z \\
\hline
Cc & t\ipafont{S} & Ii & \ipafont{I} & Oo & \ipafont{6} & Uu & \ipafont{U} & &\\
\hline
Dd & d & Jj & \ipafont{Z} & Pp & p & Vv & v & & \\
\hline
Ee & \ipafont{E} & Kk & k & Qq & \ipafont{T} & Ww & w & & \\
\hline
Ff & f & Ll & l & Rr & r & Xx & \ipafont{S} & & \\
\hline
\end{tabular}
\\\\\\
\fontsize{12pt}{12pt}\selectfont
For other sounds, the symbol ($_\prime$) can be used to modify certain symbols. For vowels it converts them into their long forms \\
(Written in \LaTeX\ as "\_\textbackslash prime")\vspace{3em}\\ % what the fuck?
%
%
\fontsize{10pt}{10pt}\selectfont List of alternate sounds:
\vspace{-1.5em}
\begin{table}[tbh!]
\begin{tabular}{||c|c||}
\specialrule{.2em}{1em}{0em}
\hline
Symbol & Sound \\
\hline
$a_\prime$ & \ipafont{eI} \\
\hline
$e_\prime$ & \ipafont{i:} \\
\hline
$i_\prime$ & \ipafont{aI} \\
\hline
$o_\prime$ & \ipafont{oU} \\
\hline
$u_\prime$ & \ipafont{u:} \\
\hline
$q_\prime$ & \ipafont{D} \\
\hline
\end{tabular}
\end{table}
\vspace{-0.5em}\\
"$q_\prime$" is very uncommon but does exist, primarily for words taken from other languages being written in Mata's alphabet. \\
For example, "The" would be written as "$q_\prime^h$". The other symbols in this will be touched on shortly.
\vspace{5em} \\
%
Mata very rarely uses the normal symbols for vowels outside of the beginning of words, or pluralisation. Instead,
Mata uses diacritics to denote when a vowel follows a consonant.\\\\
\raggedright A table of the diacritic marks follows:
\vspace{-1.5em}
\begin{table}[tbh!]
%\centering
\begin{tabularx}{21.25em}{|| >{\centering\arraybackslash}m{3em} | >{\centering\arraybackslash}m{7em} | >{\centering\arraybackslash}m{7em} ||}
\specialrule{.2em}{1em}{0em}
\hline
Symbol & Equivalent vowel & \LaTeX \\
\hline
$\hat{\circ}$ & a  & \tbs hat\{...\}\\
\hline
$\ddot{\circ}$ & e & \tbs ddot\{...\}\\
\hline
$\tilde{\circ}$ & i & \tbs tilde\{...\}\\
\hline
$\dot{\circ}$ & o & \tbs dot\{...\}\\
\hline
$\breve{\circ}$ & u & \tbs breve\{...\}\\
\hline
\end{tabularx}
\end{table}
\\ \vspace{1em}
Additionally there are two other diacritics other than for vowels. \\
%To asperate a sound, a superscript 'h' can be appended. \\
%To extend a sound, a bar can be added above or below the symbol. 
%If the bar is above another diacritic, it will extend the sound of the diacritic rather than the base consonant. \\
\begin{tabular}{||c|c||}
\specialrule{.2em}{1em}{0em}
Symbol & Sound \\\hline
$\circ^h$ & Aspirated\\\hline
$\bar{\circ}$ / $\circ$ & Long Sound\\\hline
\end{tabular}
%\typeout{Extra Row Height:}
%\showthe\extrarowheight
\setlength\extrarowheight{5pt}
\begin{tabular}{||c|c||}
\specialrule{.2em}{1em}{0em}
Symbol & Sound \\\hline
$\hat{\bar{S}}$ & Long S, short a \\\hline
$\bar{\hat{S}}$ & Short S, long a \\\hline
\end{tabular}
\setlength\extrarowheight{0pt}
\newpage

\subsection{Special Characters}
Mata uses more than a basic alphabet. There are a few special characters with dedicated meanings. All special characters from now on will have their sounds written in the Mata alphabet.
\begin{table}[tbh!]
\begin{tabular}{||M{3em}|M{7em}|M{3em}|M{7em}||}
\specialrule{.2em}{1em}{0em}
\hline
Symbol & \LaTeX & Sound & Meaning \\
\hline
$\mathcal{E}$ & \tbs mathcal\{E\} & $e_\prime$ & The speaker (I) \\
\hline
$\mho$ & \tbs mho & $\hat{w}$ & The listener (you singular) \\
\hline
$\emptyset$ & \tbs emptyset & $\hat{n}$ & 0 / nothing \\
\hline
$\xi$ & \tbs xi & $a^h$ & Everything \\
\hline
$\sigma$ & \tbs sigma & $\ddot{s}_\prime g$ & Level of significance \\
\hline
\end{tabular}
\end{table}

\subsection{Operators}
Operators are special characters which describe properties or modifications to lexical objects. In sentences, they are particles that fit between objects.
\setlength{\LTleft}{0cm}
\newcommand{\optable}{
\begin{longtable}{||M{5em}|M{11em}|M{5em}|M{12em}||}
\specialrule{.2em}{0em}{0em}
\hline
Symbol & \LaTeX & Sound & Meaning \\ \specialrule{.2em}{0em}{0em} \hline
\endfirsthead
\specialrule{.2em}{1em}{0em}
\hline
Symbol & \LaTeX & Sound & Meaning \\ \specialrule{.2em}{0em}{0em} \hline
\endhead
\hline
\endfoot
\hline
\endlastfoot} 

\newcommand{\breakoptable}{\specialrule{.2em}{0em}{0em}\hline\multicolumn{4}{?c?}{Next page...} \\ \specialrule{.2em}{0em}{0em}\hline\end{longtable}\newpage\optable}

\optable
$+$ & + & $p\breve{l}s$ & Addition \\ \hline
$-$ & - & $\tilde{m}\breve{n}s$ & Subtraction \\ \hline
$\times$ & \tbs times & $\breve{m}l\hat{t}$ & Multiplication \\ \hline
$\div$ & \tbs div & $\tilde{d}\tilde{v}j$ & Division \\ \hline
$\ni$ / $\in$ & \tbs ni / \tbs in & $\dot{k}$ or $\dot{k}_\prime$ & Contains / Is contained in \\ \hline
$\not\ni$ / $\not\in$ & \tbs not\tbs ni / \tbs not\tbs in & $\dot{n}\ni$ or $\dot{n}_\prime\ni$ & Does not contain / is not contained in \\ \hline
$\circ$ & \tbs circ & $k$ (short) & Used to access an object's child \\ \hline
$.$ & . & \emph{also} $k$ (short) & Decimal place \\ \hline
$\neg$ & \tbs neg (often on keyboards) & $\dot{n}$ or $\dot{n}_\prime$ & Not, for saying something is false \\ \hline
$\dashv$ & \tbs dashv & $\dot{m}$ or $\dot{m}_\prime$ & Not, for inverting \\ \hline
$\vdash$ & \tbs vdash & $oyl$ & Was true and is still true. \\ \hline
$\updownarrow$ & \tbs updownarrow & $vr$ & Vertical \\ \hline
$\leftrightarrow$ & \tbs leftrightarriw & $\dot{h}r$ & Horizontal \\ \hline
$=$ & = & $a^h$ & is/equals \\ \hline
$\ne$ & \tbs ne & $na^h$ & is not/not equal \\ \hline
$\overset{*}{=}$ & \tbs overset\{*\}\{=\} & $o\hat{r}$ & Requires/depends on \\ \hline
$\supset$ / $\subset$ & \tbs supset / \tbs subset & $\dot{s}$ or $\dot{s}_\prime$ & Is a subset of / contains the set... (plural of $\in$ \\ \hline
$\not\supset$ / $\not\subset$ & \tbs not\tbs supset / \tbs not\tbs subset & $\dot{n}\supset$ or $\dot{n}_\prime\supset$ & negative of subset/superset \\ \hline

\breakoptable

$?$ & ? & $v^h$ & Convert an operator to a question \\ \hline
$\{$ / $\}$ & \tbs\{ / \tbs\} & $\hat{l}$ & Open/Close a set. "\}" does not have to be pronounced at the end of a sentence. \\ \hline
$|$ & | & $\breve{X}$ & And, in a set \\ \hline
$\cap$ & \tbs cap & $\hat{X}_\prime y$ & And, between sets (intersection) \\ \hline
$\cup$ & \tbs cup & $\breve{C}y$ & Or, between sets (union) \\ \hline
$\therefore$ & \tbs therefore & $\hat{Q}\breve{r}_\prime$ & Therefore \\ \hline
$\because$ & \tbs because & $\hat{B}\breve{r}_\prime$ & Because \\ \hline
$\diamond$ & \tbs diamond & $\hat{P}$ & Is possible (can) \\ \hline
$\boxplus$ & \tbs boxplus & $\bar{\hat{S}}p\underline{l}$ & However/But (Despite this) \\ \hline
$\land$ & \tbs land & $\hat{X}$ & And, in logic \\ \hline
$\vee$ & \tbs vee & $\breve{C}$ & Or, in logic \\ \hline
$\colon$ & \tbs colon & No sound & Used optionally to separate adjacent words as Mata does not use spaces \\ \hline
\endlastfoot
\end{longtable}

\vspace{5em}
\subsection{Numbers}
Mata, like other western languages, uses base 10  and the same numerals as English (Hindu-Arabic numerals if i remember correctly). The names for numbers are as follows:\\
\begin{tabular}{||c|c||c|c||c|c||}
\specialrule{.2em}{.5em}{0em}
\hline
Numeral & Word & Numeral & Word & Numeral & Word \\ \hline
0 & $na$ & 10 & $\ddot{k}_\prime n$ & 100 & $\ddot{k}_\prime t$ \\ \hline
1 & $ak$ & 11 & $\ddot{k}_\prime nak$ & 200 & $\ddot{l}_\prime t$ \\ \hline
2 & $el$ & 20 & $\ddot{l}_\prime n$ & 300 & $\ddot{d}_\prime t$ \\ \hline
3 & $d\ddot{r}$ & 30 & $\ddot{d}_\prime n$ & 1000 & $\ddot{k}_\prime\hat{y}$ \\ \hline
4 & $f\hat{w}$ & 40 & $\ddot{f}_\prime n$ & 2000 & $\ddot{l}_\prime\hat{y}$ \\ \hline
5 & $\breve{y}n$ & 50 & $\ddot{y}_\prime n$ & 3000 & $\ddot{d}_\prime\hat{y}$ \\ \hline
6 & $\hat{\jmath}n$ & 60 & $\ddot{\jmath}_\prime n$ & & \\ \hline
7 & $\hat{p}r$ & 70 & $\ddot{p}_\prime n$ & & \\ \hline
8 & $uv$ & 80 & $\ddot{u}_\prime n$ & & \\ \hline
9 & $\ddot{y}c$ & 90 & $\ddot{c}_\prime n$ & & \\ \hline
\end{tabular}\vspace{0.25em}\\
If you wish to say a number with order above 3, you can use scientific notation with the $\hat{p}$ particle (i.e $2.5\ \text{million}\ =\ 2.5\times10^6\ =\ elk\breve{y}n\hat{p}\hat{\jmath}n$). \\
Numbers can be used for emphasis like an adjective (i.e. "very good" would be like saying "10 good"/"$\ddot{k}_\prime n:\hat{q}\breve{n}l$").\\
The ordinal variants are as follows:\\
\begin{stdtable}{15}%\begin{tabular}{||c|c|c|c|c|c|c|c|c|c||}
\hline
Number & 1 & 2 & 3 & 4 & 5 & 6 & 7 & 8 & 9 & 10 & 11 & 20 & 30 & 100 \\ \hline
Word & $a\breve{k}_\prime$ & $\ddot{l}\bar{a}$ & $\dot{d}r$ & $\breve{f}_\prime$ & $\breve{y}_\prime t$ & $\hat{\jmath}w$ & $\breve{p}\ddot{r}_\prime$ & $\hat{u}l$ & $\ddot{t}\hat{c}$ & $a\breve{k}_\prime n$ & $a\breve{k}_\prime nak$ & $a\breve{l}_\prime n$ & $a\breve{d}_\prime n$ & $a\breve{k}_\prime t$ \\ \hline
\end{stdtable}

\newpage
\section{Words}
Time for an overview of some of the words in \emph{Mata}. These lists are not extensive and at somepoint I will create a dictionary for them.
\subsection{Verbs}
\emph{Verbs} is not the best word to describe these words, they more represent the concept of the action rather than the action itself. This will be more relevant later.
Here is a list of some of the more common \emph{Verbs} (or some of the first few I created) with rough English translations:\\
%
\begin{wordtabular}{5}
To do & $\dot{d}$ & To say/speak & $\hat{S}g$ & To play & $\dot{s}nd$ & To hate & $\hat{d}\ddot{r}$ & To fish & $\hat{s}\hat{m}$ \\ \hline
To read & $k\tilde{o}c$ & To go & $\breve{m}ln$ & To consume & $v\dot{l}_\prime s$ & To get & $\dot{r}\hat{r}$ & To know & $oylr$ \\ \hline
To want & $\ddot{y}tx$ & To see & $b\hat{y}n$ & To sex & $\ddot{l}\hat{f}$ & To assume/ think/ believe & $\hat{d}rs$ & To milk & $\tilde{m}lj$ \\ \hline
To be going to/intend to (will) & $\breve{q}\hat{v}t$ & To assert (suggest is true) & $\breve{k}l\hat{t}$ & To die/kill & $ol\hat{k}$ & & & & \\ \hline
\end{wordtabular}\\
Mata supports compound \emph{verbs} as well:\\
\begin{wordtabular}[4em][4em]{5}
To sleep & $n\dot{y}c\breve{m}ln$ & To drive & $\hat{m}\tilde{x}n\breve{m}ln$ & To love & $\dashv\hat{d}\ddot{r}$ & To hunger/ thirst & $v\dot{l}_\prime s\ddot{y}tx$  & & \\\hline
\end{wordtabular}\\
In order to describe the past or future tense of a \emph{verb}, you can either
\begin{itemize}
\item Access the past or future child of the \emph{verb} (i.e. $\mathcal{E}\circ\dot{d}\circ\tilde{u}p$ is the set of things that I, the speaker, have done)
\item Apply the past or future as a function, less uses (i.e. $\tilde{u}p(\hat{s}\hat{m}$) is the state of having fished, The sentence "he is going to have fished" would be $\tilde{G}\circ\dot{d}\circ k\breve{l}\ni\tilde{u}p(\hat{s}\hat{m})$)
\end{itemize}
Here are a few example sentences:\\
\NewDocumentCommand{\sentex}{m m O{}}{\emph{#1} & #2 & \ifx#3\else(#3)\fi\\} %sentence example
\begin{tabular}{l l l}
\sentex{James says "Hello"}{$\myname\circ\hat{S}g\ni"\ddot{H}\dot{l}_\prime"$}[James'k sag koh Hello]
\sentex{James is saying "Hello"}{$\myname\circ\dot{d}\ni\hat{S}g("\ddot{H}\dot{l}_\prime")$}[James'k doh koh sag Hello]
\sentex{James is doing nothing}{$\myname\circ\dot{d}=\emptyset$}[James'k doh ah na]
\sentex{I hate you}{$\mathcal{E}\circ\hat{d}\ddot{r}\ni\mho$}[Ee'k dareh koh wa]
\sentex{I hate you a lot}{$\mathcal{E}\circ\hat{d}\ddot{r}:\ddot{k}_\prime n\ni\mho$}[Ee'k dareh keen koh wa]
\sentex{I really really love you}{$\mathcal{E}\circ\dashv\hat{d}\ddot{r}:\ddot{k}_\prime\hat{y}\ni\mho$}[Ee'k modareh keeya koh wa]
\sentex{I saw you}{$\mathcal{E}\circ b\hat{y}n\circ\tilde{u}p\ni\mho$}[Ee'k byan k'u-ip koh wa]
\sentex{I saw you reading}{$\mathcal{E}\circ b\hat{y}n\circ\tilde{u}p\ni\mho\circ\dot{d}(k\tilde{o}c)$}[Ee'k byan k'u-ip koh wa'k doh(ko-ich)]
\sentex{James can't do everything}{$\myname\circ\diamond\dot{d}\ne\xi$}[James k'padoh na ah]
\sentex{I want to go to sleep}{$\mathcal{E}\circ\ddot{y}tx\ni n\dot{y}c\breve{m}ln$}[Ee'k yetsh koh nyoch-muln]
%\sentex{I need 
\end{tabular}

\newpage
\subsection{Nouns and Pronouns}
\subsubsection{Pronouns}
\emph{Mata} supports most of the same pronouns as \emph{English} but remember that objects like pronouns are sets which can be combined together to describe different groups using the "$\cap$" and "$\cup$" symbols.\\
\begin{wordtabular}{5}
The Speaker, I & $\mathcal{E}$ & The Listener, You[sing] & $\mho$ & They/It* & $\tilde{g}$ & You[plur]$^\dagger$ & $p\hat{y}$ & They [plur]$^\ddagger$ & $g\hat{y}$ \\\hline
This/ These & $\mathcal{E}\hat{c}$ & That/ Those & $\mho\hat{c}$ & & & & & & \\\hline
\end{wordtabular}\vspace{0.25em}\\
\small{*Any physical object or person. Cannot be used for anything abstract, except for personification}\\
\small{$^\dagger$Any group of people including the listener and not including the speaker}\\
\small{$^\ddagger$Any group of people including neither the listener nor the speaker}
\subsubsection{Nouns}
\vspace{-1.5em}
\begin{wordtabular}{5}
Past & $\tilde{u}p$ & Future & $k\breve{l}$ & Name & $\tilde{n}ym$ & Night & $n\dot{y}c$ & Apple & $a\breve{p}l$ \\\hline
Car & $\hat{m}\tilde{x}n$ & Head & $\hat{l}gk$ & Mum & $\hat{m}\breve{m}_\prime$ & Dad & $\hat{p}\breve{p}_\prime$ & Child & $\tilde{k}\hat{l}n$ \\\hline
\end{wordtabular}

\newpage
\subsection{Adjectives}
Adjectives are primarily made of root words and modifier particles. For example, the word for short is equivalent to 'not vertical length'. Compound words can also be made by combining other words. For example, enjoyment is equivalent to 'doing-happiness'.\\

\centering\textbf{Roots:}\\
\begin{wordtabular}{5}
Length & $\hat{k}\hat{r}$ & Happiness & $\hat{t}_\prime\dot{r}l$ & Good & $\hat{T}\breve{n}l$ & Bad & $\breve{H}_\prime\breve{n}l$ & Empty & $\hat{v}n$ \\\hline
Area & $E\hat{r}$ & Volume & $\dot{t}\hat{r}$ & Male & $\hat{k}l$ & Female & $\hat{v}l$ & & \\\hline
\end{wordtabular}\vspace{1.25em}\\

Even though the words for length, area, and volume describe the concepts, describing an object as having the concept as a quality suggests that it has a large amount. So if I want to say that an object is long I can just say that it has length, and if i want to say that it's short I can say it has not-length. To convert these concepts into physical objects, the relevant affix would be needed. Therefore, \emph{An Area} would be $E\hat{r}\ddot{y}n$.\\\vspace{2em}

\textbf{Modifiers:}\\
These are all operators from before, however I will still reiterate the relevant ones here.\\
\begin{wordtabular}{5}
Not & $\neg$ & Vertical & $\updownarrow$ & Horizontal & $\leftrightarrow$ & & & &\\\hline
\end{wordtabular}

\newpage
\subsection{Interjections}
\subsubsection{\raggedright Frequently Used}
\vspace{-1.5em}
\begin{wordtabular}{5}
Yes & $\hat{T}$ & No & $\breve{h}_\prime$ & Thanks & $\tilde{s}k\ddot{r}l$ & Hello & $\mataHello$ & Excuse me & $i\tilde{v}\hat{n}$ \\\hline
\end{wordtabular}
\subsubsection{Expletives}
\vspace{-1.5em}
\begin{wordtabular}{5}
F*ck/ Holy Sh*t & $\hat{y}v\hat{l}\bar{r}$ & B*tch/ Idiot & $k\hat{r}n\hat{t}$ & & & & & & \\\hline
\end{wordtabular}

\end{document}